% #############################################################################
% RESUMO em Português
% !TEX root = ../main.tex
% #############################################################################
% reset acronyms
\acresetall
% use \noindent in first paragraph
\noindent Linguagens de programação baseadas em tipos de sessão e lógica linear, como CLASS, garantem segurança de memória e ausência de deadlocks por construção. A Linear Session Abstract Machine (SAM) forneceu o primeiro modelo formal para a execução sequencial de processos em linguagens à base de sessões, permitindo tempos de execução várias ordens de magnitude melhores do que interpretadores existentes que criam threads para cada processo concorrente. No entanto, a SAM, enquanto interpretador de alto nível, introduz custos de desempenho quando comparada com execuções de programas compilados. Este projeto introduz a primeira representação intermédia (IR) para linguagens baseadas em sessões, que pode ser utilizada para compilar programas tipados por sessões em código imperativo de baixo nível extremamente eficiente. Foi desenvolvido um esquema de compilação para CLASS, que traduz a linguagem de alto nível para esta IR, mantendo todas as garantias de segurança e liveness da linguagem fonte. Foi implementada uma versão deste esquema que cobre todas as funcionalidades da linguagem e que implementa várias técnicas de otimização. É agora possível compilar e executar programas escritos em CLASS, atingindo desempenhos várias ordens de magnitude melhores do que a SAM.
